\documentclass[11pt,a4paper]{article}

\usepackage{amsmath,amssymb,amsfonts,amsthm}
\usepackage{bm}
\usepackage{geometry}
\usepackage{hyperref}
\usepackage{physics}
\usepackage{enumitem}
\usepackage{longtable}
\geometry{margin=2.5cm}

\title{Einstein--\AE ther Theory: A Detailed Review}
\author{}
\date{2026}

\begin{document}

\maketitle

\begin{abstract}
Einstein--\AE ther (EA) theory is one of the most extensively studied
generally covariant modifications of General Relativity (GR) containing
a dynamically preferred local rest frame.
The theory introduces a unit timelike vector field
(the ``\AE ther'') that spontaneously breaks local Lorentz symmetry
while preserving diffeomorphism invariance.
It provides a mathematically consistent framework for studying
Lorentz violation in gravity, serves as the low-energy limit of
Ho\v rava gravity, and has become an important testing ground for
questions concerning preferred time, cosmology, black holes,
gravitational waves, dark matter phenomenology, and quantum gravity.
This review presents the foundations, mathematical structure,
experimental constraints, cosmological applications, and open problems
of Einstein--\AE ther theory.
\end{abstract}

\tableofcontents

\section{Introduction}

The standard formulation of General Relativity possesses
local Lorentz invariance.
At each spacetime point one can choose a local inertial frame,
and no physically preferred timelike direction exists.

Nevertheless several motivations suggest studying theories with
preferred frames:

\begin{enumerate}
\item Quantum gravity may introduce a fundamental time scale.
\item The cosmic microwave background defines a cosmological rest frame.
\item Canonical quantization often singles out time.
\item The problem of time in quantum gravity remains unresolved.
\item Some approaches to quantum gravity predict Lorentz violation.
\item Preferred-frame effects may mimic dark matter or dark energy.
\end{enumerate}

Einstein--\AE ther theory was developed mainly by
:contentReference[oaicite:0]{index=0}
and collaborators as a generally covariant realization of a preferred
timelike direction.

Unlike older Lorentz-violating theories, the preferred frame is not
introduced externally.
Instead it is generated dynamically by a vector field
$u^a(x)$ satisfying

\begin{equation}
u^a u_a = -1.
\end{equation}

Thus Lorentz symmetry is broken spontaneously rather than explicitly.

\section{Historical Background}

Important predecessors include:

\begin{itemize}
\item Dirac's large-number cosmology.
\item Vector-tensor theories.
\item Rosen's bimetric gravity.
\item Preferred-frame theories.
\item Lorentz-violating quantum field theories.
\end{itemize}

Modern EA theory emerged around 2000--2004 through a series of papers by
Jacobson and Mattingly.

The key innovation was to combine

\begin{enumerate}
\item diffeomorphism invariance,
\item a dynamical preferred frame,
\item second-order field equations.
\end{enumerate}

This yielded a predictive effective field theory.

\section{Fundamental Fields}

The theory contains

\begin{enumerate}
\item metric $g_{ab}$,
\item unit timelike vector field $u^a$.
\end{enumerate}

Constraint:

\begin{equation}
u^a u_a=-1.
\end{equation}

Physically:

\begin{equation}
u^a=
\text{4-velocity of the local preferred frame}.
\end{equation}

At each event the \AE ther selects a unique local time direction.

Hence the tangent space no longer possesses full Lorentz symmetry.

\section{Action}

The most general two-derivative action is

\begin{equation}
S=
\frac{1}{16\pi G}
\int d^4x
\sqrt{-g}
\left(
R + L_u
+\lambda(u^a u_a+1)
\right),
\end{equation}

where

\begin{equation}
L_u=
-K^{ab}{}_{mn}
\nabla_a u^m
\nabla_b u^n .
\end{equation}

The tensor

\begin{equation}
K^{ab}{}_{mn}
=
c_1 g^{ab}g_{mn}
+c_2 \delta^a_m\delta^b_n
+c_3 \delta^a_n\delta^b_m
-c_4 u^a u^b g_{mn}
\end{equation}

contains four dimensionless couplings

\begin{equation}
c_1,c_2,c_3,c_4.
\end{equation}

The Lagrange multiplier $\lambda$
enforces normalization.

Therefore the theory is characterized by four independent parameters.

\section{Field Equations}

Variation with respect to $g_{ab}$ yields

\begin{equation}
G_{ab}
=
T^{\rm ae}_{ab}
+
8\pi G T^{\rm matter}_{ab}.
\end{equation}

Variation with respect to $u^a$ gives

\begin{equation}
\nabla_a J^a{}_m
+
c_4 a_a\nabla_m u^a
=
\lambda u_m ,
\end{equation}

where

\begin{equation}
J^a{}_m
=
K^{ab}{}_{mn}
\nabla_b u^n.
\end{equation}

Together these define a coupled metric-vector system.

\section{Spontaneous Lorentz Symmetry Breaking}

A crucial conceptual point is that Lorentz symmetry of the laws
and Lorentz symmetry of solutions are distinct.

The action remains generally covariant.

However vacuum solutions satisfy

\begin{equation}
\langle u^a\rangle \neq 0.
\end{equation}

Thus the vacuum selects a preferred direction.

This is analogous to:

\begin{itemize}
\item ferromagnetism,
\item Higgs mechanism,
\item superfluids.
\end{itemize}

The symmetry breaking is therefore spontaneous.

This distinction is important because explicit Lorentz breaking often
creates consistency problems with energy-momentum conservation,
whereas spontaneous breaking generally does not.

\section{Mode Decomposition}

Linearization about Minkowski space yields three propagating sectors.

\subsection{Spin-2}

Ordinary graviton.

Velocity:

\begin{equation}
v_2^2
=
\frac{1}{1-c_{13}}
\end{equation}

where

\begin{equation}
c_{13}=c_1+c_3.
\end{equation}

\subsection{Spin-1}

Vector \AE ther mode.

Velocity:

\begin{equation}
v_1^2
=
\frac{
c_1-\frac12 c_1^2+\frac12 c_3^2
}
{
c_{14}(1-c_{13})
}.
\end{equation}

\subsection{Spin-0}

Scalar mode.

Velocity:

\begin{equation}
v_0^2
=
\frac{
c_{123}(2-c_{14})
}{
c_{14}(1-c_{13})(2+c_{13}+3c_2)
}
\end{equation}

where

\begin{equation}
c_{123}=c_1+c_2+c_3.
\end{equation}

Unlike GR, three different gravitational wave sectors may exist.

\section{Relation to Effective Metrics}

Each mode propagates on a different effective geometry.

Schematically,

\begin{equation}
g^{(i)}_{ab}
=
g_{ab}
+
\alpha_i u_a u_b .
\end{equation}

This is closely related to the ``disformal'' metric

\begin{equation}
\tilde g_{ab}
=
g_{ab}
-
2u_a u_b .
\end{equation}

which you mentioned in previous discussions.

Such metrics naturally arise in the characteristic structure of the
field equations.

Consequently the theory can be viewed as possessing several causal cones.

The physical light cone of matter need not coincide with
the propagation cones of all gravitational modes.

\section{PPN Formalism}

The Parametrized Post-Newtonian framework provides precision tests.

The most important parameters are

\begin{equation}
\alpha_1,
\qquad
\alpha_2,
\end{equation}

which measure preferred-frame effects.

Solar-system experiments require

\begin{equation}
|\alpha_1|
\ll 1,
\qquad
|\alpha_2|
\ll 1.
\end{equation}

These constraints restrict the allowed combinations of
$c_i$.

Remarkably a large parameter region survives.

\section{Gravitational Waves}

The observation of the
:contentReference[oaicite:1]{index=1}
binary neutron-star merger dramatically constrained modified gravity.

The arrival times of photons and gravitational waves agree to
approximately one part in $10^{15}$.

Consequently

\begin{equation}
v_2 \approx c.
\end{equation}

This eliminates much of parameter space.

Current viable EA models lie near

\begin{equation}
c_{13}\approx 0.
\end{equation}

\section{Cosmology}

For FLRW spacetime

\begin{equation}
ds^2
=
-dt^2
+a(t)^2 d\vec{x}^2,
\end{equation}

the \AE ther aligns naturally with cosmic time:

\begin{equation}
u^a=(1,0,0,0).
\end{equation}

Modified Friedmann equations become

\begin{equation}
H^2
=
\frac{8\pi G_{\rm cos}}{3}\rho .
\end{equation}

Effective Newton constants differ:

\begin{equation}
G_{\rm cos}
\neq
G_{\rm local}.
\end{equation}

This affects

\begin{itemize}
\item nucleosynthesis,
\item CMB anisotropies,
\item structure formation.
\end{itemize}

\section{Connection to the CMB Frame}

An interesting conceptual observation is that the CMB defines a
distinguished cosmological rest frame.

In standard GR this is merely a property of a particular solution.

In Einstein--\AE ther theory a preferred timelike direction exists
already at the level of dynamical fields.

Thus:

\begin{center}
CMB frame $\neq$ \AE ther field.
\end{center}

However cosmological solutions often align them.

This is exactly the distinction between

\begin{itemize}
\item preferred state (GR),
\item preferred law (EA).
\end{itemize}

\section{Black Holes}

Black holes exhibit novel features.

Different modes possess different horizons.

One finds:

\begin{enumerate}
\item metric horizon,
\item spin-1 horizon,
\item spin-0 horizon.
\end{enumerate}

Most remarkable is the existence of a

\begin{equation}
\text{universal horizon}.
\end{equation}

This hypersurface traps even infinitely fast signals.

It therefore generalizes the notion of a black-hole horizon
for Lorentz-violating theories.

\section{Relation to Ho\v rava Gravity}

The strongest modern motivation comes from
:contentReference[oaicite:2]{index=2}
gravity.

Ho\v rava gravity assumes a preferred foliation

\begin{equation}
t={\rm const}.
\end{equation}

At low energies it becomes equivalent to a special sector of
Einstein--\AE ther theory.

Specifically

\begin{equation}
u_a
=
\frac{\nabla_a T}
{\sqrt{-\nabla_bT\nabla^bT}}
\end{equation}

for a scalar ``khronon'' field $T$.

This is called hypersurface-orthogonal Einstein--\AE ther theory.

\section{Dark Matter Connections}

Several mechanisms have been proposed.

\subsection{Modified Dynamics}

The \AE ther contributes additional gravitational forces.

Some galactic phenomena can be reproduced.

\subsection{MOND-like Extensions}

Einstein-\AE ther modifications can generate

\begin{equation}
a \sim \sqrt{a_0 g_N}
\end{equation}

behavior.

\subsection{Dark Sector Field}

The \AE ther itself behaves as a cosmic medium.

Nevertheless no fully successful replacement for
cold dark matter has emerged.

Current consensus remains that EA theory does not eliminate
the need for dark matter.

\section{Relation to Global Time}

This topic is particularly relevant to
Stueckelberg--Horwitz--Piron (SHP) ideas.

EA theory introduces a preferred local timelike vector field.

However it does \emph{not} provide a universal global evolution
parameter analogous to

\begin{equation}
\tau
\end{equation}

of SHP.

The preferred time is local and dynamical.

Therefore:

\begin{center}
SHP global time
$\neq$
Einstein--\AE ther preferred frame.
\end{center}

Nevertheless both approaches share the intuition that spacetime
may possess more temporal structure than standard relativistic
kinematics suggests.

\section{Current Experimental Status}

Surviving theories must satisfy:

\begin{enumerate}
\item Solar-system constraints.
\item Binary pulsar constraints.
\item Gravitational-wave constraints.
\item Cosmological constraints.
\item Stability requirements.
\item Absence of vacuum Cherenkov radiation.
\end{enumerate}

After GW170817 the viable parameter space became small but remains
non-empty.

Thus Einstein--\AE ther theory is not excluded.

\section{Open Problems}

Major unresolved issues include:

\begin{enumerate}
\item Quantum formulation.
\item Renormalization.
\item Relation to fundamental quantum gravity.
\item Black-hole thermodynamics.
\item Dark matter phenomenology.
\item Singularities and cosmic censorship.
\item Relation to the problem of time.
\item Emergence of Lorentz symmetry.
\end{enumerate}

\section{Assessment}

Einstein--\AE ther theory occupies a unique position among modified
gravity theories.

Advantages:

\begin{itemize}
\item Generally covariant.
\item Dynamical preferred frame.
\item Well-posed field equations.
\item Rich phenomenology.
\item Connection to Ho\v rava gravity.
\item Strong observational testability.
\end{itemize}

Disadvantages:

\begin{itemize}
\item Additional parameters.
\item No compelling quantum completion.
\item Strong observational constraints.
\item No established solution to dark matter.
\end{itemize}

From a conceptual standpoint, the theory demonstrates that
Lorentz symmetry and general covariance are logically distinct.
One may preserve diffeomorphism invariance while introducing a
physically meaningful preferred timelike direction.

For researchers interested in preferred time,
cosmological frames, SHP theory, or alternatives to strict local
Lorentz invariance, Einstein--\AE ther theory remains one of the
most mathematically developed and experimentally constrained
frameworks available.

\begin{thebibliography}{99}

\bibitem{Jacobson2001}
T. Jacobson and D. Mattingly,
``Gravity with a Dynamical Preferred Frame.''

\bibitem{Jacobson2008}
T. Jacobson,
``Einstein-Aether Gravity: A Status Report.''

\bibitem{ElingJacobson}
C. Eling, T. Jacobson, and D. Mattingly,
``Einstein-Aether Theory.''

\bibitem{JacobsonHorava}
T. Jacobson,
``Extended Horava Gravity and Einstein-Aether Theory.''

\bibitem{Foster}
B. Foster,
``Strong Field Effects on Binary Systems in Einstein-Aether Theory.''

\end{thebibliography}

\end{document}
